\subsection{RQ1: control boundary rather than product label}
The temporal category contains exclusive execution, queueing, context switch, kernel-priority and subgraph policies. Clockwork's predictable worker design and REEF's priority/preemption design share a latency objective but control different execution units \cite{clockwork,reef}. PipeSwitch's context switch latency concerns a transition mechanism; it is not automatically a request p99 outcome \cite{pipeswitch}. SPLIT changes the model graph and scheduling blocks, and SwitchFlow interleaves inference subgraphs \cite{split,switchflow}. Grouping them only as ``time-slicing'' would hide the preemption and memory requirements that determine whether a tail bound transfers to another service.

The MPS family controls process concurrency and active-thread allocation. GPUlet, GSLICE, iGniter and DeInfer add provisioning, admission or interference modeling around those primitives \cite{gpulet,gslice,igniter,deinfer}. MIG divides a supported GPU into predefined instance layouts; GNN-Triton exposes slice fit, whereas Dyna-P and MIGServe add dynamic placement or migration \cite{gnntriton,dynap,migserve}. Hybrid systems such as BOER add MPS inside MIG, and a second software controller may change the outcome even when the hardware profile is held fixed \cite{boer}. API systems have still another control point: Torpor measures transfer paths, while Tally and Orion use intercepted calls to regulate execution \cite{torpor,tally,orion}. These categories are therefore a map of what can be controlled and measured, not a guarantee of equivalent implementations.

\subsection{RQ2: outcome and baseline matrix}
\begin{longtable}{P{0.21\linewidth}P{0.34\linewidth}P{0.34\linewidth}}
\caption{Why apparently similar outcomes cannot be pooled without a matched denominator.}\\
\toprule Outcome & Example of direct measurement & Comparison boundary\\\midrule\endfirsthead
\toprule Outcome & Example of direct measurement & Comparison boundary\\\midrule\endhead
SLO goodput & Clockwork measures requests that satisfy the deadline; InferFair reports local goodput and a slowdown-fairness ratio \cite{clockwork,inferfair}. & Requires the same SLO and arrivals; neither is raw request throughput or an adversarial property.\\
Latency percentile & REEF and XSched use p99; Torpor reports a p98 tail; SoCC interference uses P90 TBT \cite{reef,xsched,torpor,interference_deeper}. & Percentile, request/token unit, load and baseline must align. A mean or TBT cannot substitute for p99 request latency.\\
Capacity and density & GSLICE uses GUE and iGniter counts instances/violations; MIGServe reports sustainable server rate at 90\% SLO attainment \cite{gslice,igniter,migserve}. & Efficiency ratio, number of GPUs and requests/s under an attainment threshold are different denominators.\\
Job time and energy & Dyna-P reports JCT/makespan; EnerTune measures pairing energy and latency; WattWiser evaluates whole-server power \cite{dynap,enertune,wattwiser}. & Completion time is not online tail latency, and server power cannot establish isolated one-slice energy.\\
Isolation & KRYPTON includes a specific propagation experiment; DeInfer and SoCC characterize performance interference \cite{krypton,deinfer,interference_deeper}. & A process boundary or SM quota does not by itself imply security or fault containment.\\
\bottomrule\end{longtable}

Metric choice can reverse an apparent ranking. A controller may complete more total requests while fewer meet a hard SLO. A profiler that raises occupancy may introduce head-of-line blocking under bursty arrivals. Cross-study normalized speedups are especially fragile when one paper uses a solo-model reference and another uses a static-partition or adaptive scheduling baseline. The extraction register retains the original comparator and statistic, allowing a reader to reconstruct each bounded statement.

\subsection{RQ3: hardware and workload strata}
Server-class NVIDIA GPUs with MIG support are not equivalent to an RTX model lacking MIG or to AMD CU masks. Dyna-P's A30 profile result and GNN-Triton's MIG slice capacity inform partitioning, while DeInfer's RTX 3080 MPS experiment informs software sharing without a MIG comparison \cite{dynap,gnntriton,deinfer}. VastPipe reports A100 MPS and a distinct MI100 CU-mask stratum; treating them as repeated measurements of one product would misstate its evaluation \cite{vastpipe}. Edge studies on Jetson or other small accelerators contribute scheduling ideas but change memory capacity, supported partition controls and launch overhead \cite{miriam,split,edge_sky}. The evidence for a given GPU family should therefore be presented with the exact SKU and runtime when available, and otherwise as an indirect hypothesis.

Model memory footprint and batch size alter overlap. GNN-Triton documents a local slice thrashing point as more models are loaded \cite{gnntriton}; a smaller MIG instance can raise simultaneous placement count but cause capacity pressure. ITIF shares transformer computation itself and cannot transfer its reported performance ratio to unrelated model families \cite{itif}. Deft's optimal synchronization frequency varies with image or batch size \cite{deft}. FaST-GShare's eight RNNT pods and their quota are a specific concurrency configuration \cite{fastgshare}. LLM decode introduces variable token lengths and sustained memory traffic: the SoCC co-tenant study directly measures P90 TBT rather than merely end-to-end average latency \cite{interference_deeper}. These are reasons to stratify by architecture and request shape before concluding that a mechanism ``improves utilization.''

Arrival and tenancy also matter. A solo baseline, a pair of fixed-rate models, a bursty Poisson trace, and a multi-stage serverless workflow pose different scheduling problems. iGniter's four vision-model constant-arrival deployment and its initial p99 recovery transient delimit the result \cite{igniter}. AdMix distinguishes a local single-GPU LLM stratum from a two-GPU tensor-parallel stratum \cite{admix}. FluidFaaS and ESG combine pipeline or workflow placement with local sharing; their end-to-end SLO percentages are relevant to deployed systems but do not isolate a MIG or MPS effect \cite{fluidfaas,esg}. A final decision should repeat the comparison under a representative trace and record both steady-state and transition measurements.
